\begin{tikzpicture}[
  font=\scriptsize,
  execute at begin node={\hyphenpenalty=10000\exhyphenpenalty=10000},
  tujuan/.style={draw, rounded corners, fill=praditagreen!15, text width=11.6cm, minimum height=0.75cm, align=center},
  tahap/.style={draw, rounded corners, fill=blue!8, text width=1.8cm, inner xsep=1.5pt, minimum height=1.5cm, align=center},
  tekan/.style={tahap, fill=orange!12},
  kapabilitas/.style={tahap, minimum height=1.35cm},
  kandidat/.style={tahap, minimum height=1.35cm},
  lanjut/.style={tahap, minimum height=1.35cm, dashed},
  prasyarat/.style={draw, rounded corners, fill=orange!12, text width=11.6cm, minimum height=0.6cm, align=center},
  lain/.style={draw, dashed, rounded corners, text width=11.6cm, minimum height=0.55cm, align=center},
  ket/.style={font=\tiny, anchor=west},
  panah/.style={-{Stealth[length=2mm]}, thick}
]
  % Satu baris per tingkat penurunan, supaya kolom dapat dibaca dari tujuan sampai kandidat.
  \node[tujuan] (tujuan) at (5.0,0) {\textbf{Tujuan strategis:} Satu Pasien, Satu Rekam Medis. Layanan terintegrasi lintas fasilitas dan ekosistem mitra};

  \node[ket] at (-0.98,-0.74) {\textit{Value stream} ``Pasien menerima layanan rawat jalan'' dan keluhan 2025 (IT40513 Bab 5)};
  \node[tahap] (t1) at (0.0,-1.75) {1. Memesan layanan\\{\tiny\mdseries 180 keluhan}};
  \node[tekan] (t2) at (2.0,-1.75) {2. Hadir dan terverifikasi\\{\tiny\mdseries 640 keluhan}};
  \node[tekan] (t3) at (4.0,-1.75) {3. Diperiksa dan didiagnosis\\{\tiny\mdseries 720 keluhan}};
  \node[tahap] (t4) at (6.0,-1.75) {4. Menerima obat\\{\tiny\mdseries 350 keluhan}};
  \node[tahap] (t5) at (8.0,-1.75) {5. Menyelesaikan pembayaran\\{\tiny\mdseries 190 keluhan}};
  \node[tahap] (t6) at (10.0,-1.75) {6. Ditindaklanjuti\\{\tiny\mdseries 70 keluhan}};
  \draw[panah] (tujuan.south -| t5.north) -- node[right, font=\tiny] {mengarahkan} (t5.north);

  \node[kapabilitas] (k1) at (0.0,-3.5) {CAP-L03 Penjadwalan Layanan};
  \node[kapabilitas] (k2) at (2.0,-3.5) {CAP-P01 Klaim dan Penjaminan};
  \node[kapabilitas] (k3) at (4.0,-3.5) {CAP-L06 Farmasi dan Laboratorium};
  \node[kapabilitas] (k4) at (6.0,-3.5) {CAP-L06 Farmasi dan Laboratorium};
  \node[kapabilitas] (k5) at (8.0,-3.5) {CAP-P01, CAP-P04 Klaim, Keuangan};
  \node[kapabilitas] (k6) at (10.0,-3.5) {CAP-L03 Penjadwalan Layanan};

  \node[kandidat] (a1) at (0.0,-5.2) {API-01 Jadwal dan Janji Temu};
  \node[lanjut] (a2) at (2.0,-5.2) {Verifikasi penjaminan\\{\tiny\mdseries kandidat lanjutan}};
  \node[kandidat] (a3) at (4.0,-5.2) {API-02 Hasil Pemeriksaan};
  \node[kandidat] (a4) at (6.0,-5.2) {API-08 Resep Apotek};
  \node[kandidat] (a5) at (8.0,-5.2) {API-03 Klaim, API-07 Mutasi Bank};
  \node[kandidat] (a6) at (10.0,-5.2) {API-01, lewat pengingat kontrol};

  \foreach \kolom in {1,...,6} {
    \draw[panah] (t\kolom.south) -- (k\kolom.north);
    \draw[panah] (k\kolom.south) -- (a\kolom.north);
  }

  \node[prasyarat] (prasyarat) at (5.0,-6.4) {\textbf{API-06 Identitas dan Akses} (CAP-P02): prasyarat lintas tahap, tidak muncul pada satu tahap saja};
  \node[lain] at (5.0,-7.15) {Penerima nilai lain, di luar \textit{value stream} ini: API-04 Rekap Perusahaan, API-05 Rekap Klinik};
\end{tikzpicture}
